\documentclass[12pt]{amsart}


\setlength{\textheight}{23cm} 
\setlength{\textwidth}{17.4cm} 
\setlength{\topmargin}{-1.3cm} 
\addtolength{\oddsidemargin}{-2cm} 
\addtolength{\evensidemargin}{-2cm} 

%\input s:/acentos.tex% fazer acentuação como no word

\newcommand{\ai} {\vbox to 7pt{\hbox to 7pt{\vrule height 7pt width 7pt}}}

%\newcommand{\overrightarrow} {\overrightarrow}



\usepackage{amssymb}

\def\R{\mathbb R}
\def\C{\mathbb C}
\def\Z{\mathbb Z}
\def\N{\mathbb N}
\def\Q{\mathbb Q}


\def\pn{\par\noindent}
\def\dps{\displaystyle}
\def\vs{\vskip}
\def\hs{\hskip}
\def\st {\stackrel}
\def\cn{\centerline}
\usepackage[all
%,dvips
]{xy} \xyoption{poly}
\usepackage{epsfig}

\begin{document}



\cn{\bf \large  Midterm - Representations of Quivers}

\medskip

\cn{\bf April 13th, 2026}

\bigskip

\medskip

\noindent {\it Choose up to five problems from this list to submit.}


\medskip

\begin{enumerate}

\item [1.] Let $Q$ be the quiver $\xymatrix{1&&2\ar[rr]\ar[ll]&&3}$ and let $M$ be a representation 
\bigskip
$$
\begin{array}{c}
 \xymatrix{k&&&k^3\ar[rrr]^{\left[\begin{array}{ccc} 1&0&0\\0&1&0\end{array}\right]}\ar[lll]_{\left[\begin{array}{ccc} 1&0&0\end{array}\right]}&&&k^2}.\\
\end{array}
$$


\noindent Write $M$ as a direct sum of simple representations, indecomposable projective and injective representations.




\bigskip

\item [2.] Let $Q$ be the quiver 

\
$$
\xymatrix{1&&2\ar@<0.5ex>[ll]^{\beta}\ar@<-0.5ex>[ll]_{\alpha}}
$$

\noindent For any $\lambda\in k\cup\{\infty\}$, define $M_{\lambda}$ to be the representation:

\
$$
\xymatrix{1&&2\ar@<0.5ex>[ll]^{\lambda}\ar@<-0.5ex>[ll]_1}\quad \lambda\in k;
$$
\
$$
\xymatrix{1&&2\ar@<0.5ex>[ll]^0\ar@<-0.5ex>[ll]_1}\quad \lambda=\infty.
$$

\

\begin{enumerate}
\item Show that each $M_{\lambda}$ is indecomposable.
\item Show that $M_{\lambda}\simeq M_{\mu}$ if and only if $\lambda=\mu$. In particular,
 the number of indecomposable representations depends on the choice of the field $k$.
\item Show that $\operatorname{Hom}(M_{\lambda},M_{\mu})=0$ if $\lambda\neq \mu$.
\end{enumerate}
\bigskip


\bigskip


\item[3.] Let $Q$ be a quiver, $\xymatrix{L\ar[r]^f & M\ar[r]^g&N}$ be an exact sequence in $rep\, Q$ and $X$ be representation of $Q$. 

\noindent Show that 
$\operatorname{Ker}g_*\subseteq \operatorname{Im} f_*$, where $f_*$ and $g_*$ are defined in
$$
 \xymatrix{\operatorname{Hom}(X,L)\ar[r]^{f_*}&\operatorname{Hom}(X,M)\ar[r]^{g_*} &\operatorname{Hom}(X,N)}.
$$

\bigskip

\medskip

\item[4.] Let $g: L\to M$ be an injective morphism between representations of $Q$, and let $I(i)$ be the injective
representation at vertex $i$. Show that  the map 
$$
g^*:\operatorname{Hom}(M,I(i)) \to \operatorname{Hom}(L,I(i))
$$
\noindent is surjective.

\bigskip

\medskip

\item[5.] Prove that any  projective representation corresponding to vertex $i$, $P(i)$ is indecomposable.

\bigskip

\medskip

\item[6.] Compute a projective resolutions for representations $S(1)$
and $S(2)$ for quiver

\bigskip

$$
Q: \qquad \xymatrix{1\ar[rr]\ar@/^0.9pc/[rrrr]&&2\ar[d]\ar[rr]&&3\ar[lld]\\ &&4&&}
$$


\bigskip

\medskip


\item[7.] Let $M=(M_i,\psi_{\alpha})$ be a representation of $Q$. Prove that for any vertex $i$ in $Q$, there is an isomorphism of vector
spaces: 

$$\operatorname{Hom}(M,I(i))\simeq M_i.$$

\medskip

\bigskip


\item[8.] Let $i$ and $j$ be vertices in the quiver $Q$. Prove that the vector space $\operatorname{Hom}(I(i),I(j))$
has a basis consisting of all paths from $j$ to $i$ in $Q$. In particular, $\operatorname{End}(I(i))\simeq k$.


\end{enumerate}

\end{document}
